% ============================================================
%  Part VIII — Word order
% ============================================================
\gpart{VIII}{Word Order}

% ------------------------------------------------------------
\gsec{The verb-second rule}

\gintro{The single most reliable rule in German syntax: in a main clause the conjugated verb occupies the second position, whatever comes first.}

In a German main clause the \textbf{conjugated verb is the second element}. Not
the second word — the second unit of meaning.

\begin{gtbl}{colspec={X[1.1,l] X[0.65,c] X[1.25,l]}}
  Position 1 & Verb & The rest \\
  Ich          & fahre & morgen nach Berlin. \\
  Morgen       & fahre & ich nach Berlin. \\
  Nach Berlin  & fahre & ich morgen. \\
  Weil es billig ist, & fahre & ich mit dem Bus. \\
\end{gtbl}

\begin{gnote}
Whatever you put in position 1 gets the emphasis, and the subject then slides in
straight after the verb. A whole subordinate clause counts as a single element,
which is why \textit{fahre} still comes second in the last example.
\end{gnote}

\tcap{The sentence bracket}

Everything verbal that is not the conjugated verb is parked at the end.

\begin{gtbl}{colspec={X[0.8,l] X[0.6,c] X[0.85,l] X[0.85,l]}}
  Pos.\,1 & Verb & Middle field & End \\
  Ich & habe & gestern ein Buch & gekauft. \\
  Ich & will & morgen früh & aufstehen. \\
  Er  & steht & jeden Tag um sechs & auf. \\
  Das Haus & wird & nächstes Jahr & gebaut. \\
\end{gtbl}

\tcap{Verb at the very end}

\begin{flattbl}{colspec={X[1,l] X[1.8,l]}}
  After a subordinating conjunction & \dots, weil ich müde \textbf{bin}. \\
  In a relative clause & Der Mann, der dort \textbf{steht}\dots \\
  In an indirect question & Ich weiß nicht, wo er \textbf{wohnt}. \\
\end{flattbl}

% ------------------------------------------------------------
\gsec{Order in the middle field}

\gintro{Everything between the conjugated verb and the clause-final verb is the middle field. Its internal order follows conventions rather than strict rules.}

\tcap{TeKaMoLo}

\begin{gtbl}{colspec={X[0.75,l] X[0.85,l] X[1.4,l]}}
  Slot & Asks & Example \\
  \textbf{Te}mporal & wann?  & heute, um acht, letzte Woche \\
  \textbf{Ka}usal   & warum? & wegen des Wetters, deshalb \\
  \textbf{Mo}dal    & wie?   & mit dem Bus, gern, schnell \\
  \textbf{Lo}kal    & wo? wohin? & in die Stadt, nach Hause \\
\end{gtbl}

\begin{flattbl}{colspec={X[1,l]}}
  Ich fahre \textbf{heute} \textbf{wegen des Wetters} \textbf{mit dem Bus} \textbf{in die Stadt}. \\
\end{flattbl}

\tcap{Objects}

\begin{gtbl}{colspec={X[1.2,l] X[1.8,l]}}
  Rule & Example \\
  Dativ before Akkusativ \eng{two nouns} & Ich gebe \textbf{dem Kind den Ball}. \\
  Pronoun before noun & Ich gebe \textbf{ihm} den Ball. \\
  Akkusativ before Dativ \eng{two pronouns} & Ich gebe \textbf{ihn ihm}. \\
\end{gtbl}

\begin{gnote}
The single rule underneath all three lines: \textbf{pronouns come early, and
known information comes before new information}. Two pronouns is the only case
where accusative jumps ahead of dative.
\end{gnote}

\tcap{Negation}

\begin{flattbl}{colspec={X[0.85,l] X[2.15,l]}}
  kein  & negates a noun with \textit{ein} or no article: \textit{Ich habe \textbf{kein} Auto.} \\
  nicht & negates everything else \\
  nicht & goes \textbf{before} the word it targets: \textit{nicht heute}, \textit{nicht schön} \\
  nicht & goes \textbf{late} when it negates the whole sentence, but still before the verb at the end \\
\end{flattbl}

% ------------------------------------------------------------
\gsec{Conjunctions}

\gintro{German conjunctions fall into three families, each with a different effect on where the verb sits.}

\tcap{Coordinating — position 0, word order untouched}

\begin{gtbl}{colspec={X[0.75,l] X[0.9,l] X[1.35,l]}}
  Word & Meaning & Example \\
  und     & and       & Ich komme \textbf{und} du bleibst. \\
  aber    & but       & Ich will, \textbf{aber} ich kann nicht. \\
  oder    & or        & Wir gehen \textbf{oder} wir bleiben. \\
  denn    & because   & Ich bleibe, \textbf{denn} ich bin müde. \\
  sondern & but rather & Nicht heute, \textbf{sondern} morgen. \\
\end{gtbl}

\tcap{Subordinating — verb to the end}

\begin{gtbl}{colspec={X[0.8,l] X[0.85,l] X[0.8,l] X[0.85,l]}}
  weil & because & obwohl & although \\
  dass & that & damit & so that \\
  wenn & if, whenever & bevor & before \\
  als & when \eng{one past event} & nachdem & after \\
  ob & whether & während & while \\
  falls & in case & seit(dem) & since \\
  bis & until & sobald & as soon as \\
\end{gtbl}

\tcap{Adverbial — they occupy position 1, so the verb follows}

\begin{gtbl}{colspec={X[0.8,l] X[0.85,l] X[1.35,l]}}
  Word & Meaning & Example \\
  deshalb   & therefore & \textbf{Deshalb} \textbf{bleibe} ich zu Hause. \\
  deswegen  & for that reason & \textbf{Deswegen} \textbf{kam} er nicht. \\
  trotzdem  & nevertheless & \textbf{Trotzdem} \textbf{gehe} ich hin. \\
  dann      & then & \textbf{Dann} \textbf{essen} wir. \\
  außerdem  & besides & \textbf{Außerdem} \textbf{ist} es teuer. \\
\end{gtbl}

\begin{gnote}
\textit{denn} and \textit{weil} both mean \emph{because} but behave in opposite
ways: \textit{Ich bleibe, \textbf{denn} ich \textbf{bin} müde} versus
\textit{Ich bleibe, \textbf{weil} ich müde \textbf{bin}}.
\end{gnote}


% ------------------------------------------------------------
\gsec{Infinitive clauses with zu}

\gintro{Where English uses \emph{to do something}, German uses \textit{zu} plus
an infinitive at the end of the clause. Several conjunctions build on the same
pattern.}

\begin{gtbl}{colspec={X[0.9,l] X[0.95,l] X[1.15,l]}}
  Pattern & Meaning & Example \\
  zu + Inf.       & to do          & Ich versuche, früh \textbf{zu kommen}. \\
  um \dots\ zu    & in order to    & Ich lerne, \textbf{um} zu bestehen. \\
  ohne \dots\ zu  & without doing  & Er ging, \textbf{ohne} zu zahlen. \\
  (an)statt \dots\ zu & instead of doing & Sie schlief, \textbf{statt} zu lernen. \\
\end{gtbl}

\begin{gnote}
With a separable verb the \textit{zu} goes \textbf{inside} the word, between the
prefix and the stem: \textit{auf\textbf{zu}stehen}, \textit{ein\textbf{zu}kaufen},
\textit{mit\textbf{zu}kommen}.
\end{gnote}

\tcap{No zu after these}

\begin{flattbl}{colspec={X[0.95,l] X[2.05,l]}}
  The six modals & Ich muss \textbf{gehen}. \\
  werden \eng{future} & Ich werde \textbf{gehen}. \\
  lassen, sehen, hören & Ich lasse ihn \textbf{gehen}. \textbullet{} Ich höre sie \textbf{singen}. \\
  gehen, fahren, bleiben & Ich gehe \textbf{einkaufen}. \\
\end{flattbl}

\tcap{brauchen and the double infinitive}

\begin{flattbl}{colspec={X[1,l] X[2,l]}}
  brauchen + zu & only in the negative: \textit{Du brauchst nicht \textbf{zu} kommen} \eng{you needn't come} \\
  Double infinitive & in the perfect, a modal keeps its infinitive form: \textit{Ich habe gehen \textbf{müssen}} \\
  Word order & with a double infinitive the auxiliary moves in front: \textit{\dots, weil ich \textbf{habe} gehen müssen} \\
\end{flattbl}

\begin{gnote}
\textit{um \dots\ zu} works only when both halves share a subject. If they do
not, use \textit{damit}: \textit{Ich erkläre es, \textbf{damit} du es
verstehst.}
\end{gnote}

% ------------------------------------------------------------
\gsec{Commas}

\gintro{German comma rules are grammatical rather than rhythmic. A comma marks a
clause boundary, so it is largely predictable once you can spot a subordinate
clause.}

\tcap{A comma is required}

\begin{gtbl}{colspec={X[1,l] X[2,l]}}
  Where & Example \\
  Before a subordinate clause & Ich bleibe\textbf{,} weil es regnet. \\
  Around a relative clause & Der Mann\textbf{,} der dort steht\textbf{,} ist mein Bruder. \\
  Between two main clauses & Ich komme\textbf{,} und sie bleibt. \eng{optional with und} \\
  Before aber, sondern, denn & Ich will\textbf{,} aber ich kann nicht. \\
  In a list & Äpfel\textbf{,} Birnen und Bananen \\
\end{gtbl}

\tcap{No comma}

\begin{flattbl}{colspec={X[1,l] X[2,l]}}
  Before \textit{und} joining two verbs with one subject & Ich stehe auf und gehe. \\
  Before \textit{wie} in a comparison & Er ist so groß wie ich. \\
  After a fronted element & Morgen fahre ich weg. \\
\end{flattbl}

\begin{gnote}
The comma before a \textit{zu} infinitive may be left out only when the
infinitive is bare — nothing depends on it and no conjunction precedes it:
\textit{Iris versprach mitzumachen.} As soon as the infinitive is extended, and
always after \textit{um}, \textit{ohne}, \textit{statt}, \textit{anstatt},
\textit{als} or \textit{außer}, the comma is compulsory.
\end{gnote}

\begin{gnote}
The list comma before \textit{und} — the Oxford comma — is not used in German.
\textit{Äpfel, Birnen und Bananen} never takes a comma before \textit{und}.
\end{gnote}
